% p3-04-pellis-expansion

\section{P3 §4. The Additive Pellis Expansion}

Figure (fig-p3-ch4-pellis). Pellis expansion triptych (Paper III) --- greedy phi-algorithm /
                                         interference / convergence funnel.

  4.1 Construction via Greedy $\varphi$-Algorithm
  The Pellis greedy expansion follows the same logical pattern as the Rényi $\beta$-expansion
  for base beta = $\varphi$, but with relaxed digit constraints and additional hierarchical
  structure.

  Algorithm 4.1 (Pellis greedy expansion to depth D).

  Input: x > 0, depth parameter D $\geq$ 0.

  Initialise: r0 = x, coefficients list C = [].

  For k = 0, 1, …, D:

  1. Compute the tentative coefficient ck = rk $\cdot$ $\varphi$k.
  2. Set ck = lfloor ck rfloor (the floor; optionally the nearest integer for centred
  expansions).
  3. Set rk+1 = rk - ck $\varphi$-k.
  4. Append ck to C.

  Output: coefficient sequence (c0, c1, …, cD) and remainder RD = rD+1.

  The expansion is then

  x = sumk=0D ck $\varphi$-k + RD(x). 1

        Proposition 4.2 (Non-negativity and bound). The greedy Pellis expansion
        satisfies 0 $\leq$ RD(x) < $\varphi$-(D+1) for all x $\geq$ 0. The coefficients ck are non-negative
        integers whenever x is a positive real.

        Proof. By construction, rk+1 = rk - lfloor rk $\varphi$k rfloor $\varphi$-k $\in$ [0, $\varphi$-k). Iterating
        gives rD+1 $\in$ [0, $\varphi$-(D+1)). □

  Comparison with Zeckendorf's representation. In the standard Zeckendorf setting, the
  input is a positive integer N, the digit set is {0, 1}, and the constraint of no two
  consecutive 1s ensures uniqueness. The Pellis expansion relaxes the integer input to
  arbitrary x $\in$ R>0, permits coefficients greater than 1, and adds the hierarchical
  interference layer (Section 4.2). It reduces to a variant of the Zeckendorf
  representation when x $\in$ Z>0 and when the coefficient set is restricted.

  Comparison with Parry--Rényi $\beta$-expansions. The Rényi (1957) / Parry (1960) theory of
  $\beta$-expansions in base beta = $\varphi$ uses digits in {0, 1} with the "no two consecutive 1s"
  constraint (since $\varphi$2 = $\varphi$ + 1). The Boyd beta-expansions framework further
  characterises when such expansions are periodic (Salem/Pisot numbers). The Pellis
  expansion departs from this setup in two ways: (i) it allows arbitrary non-negative
  integer coefficients, not just binary digits; (ii) it adds an interference correction layer
  (Section 4.2) that encodes cancellations between adjacent levels. The Pellis expansion
  should be understood as a signed generalisation of the Rényi $\beta$-expansion, not a special
  case.

  4.2 Hierarchical Interference Structure
  The hierarchical Pellis expansion adds a second level of structure by decomposing the
  coefficients ck into sums involving lower-level coefficients, encoding interference
  between adjacent depth levels.

        Definition 4.3 (Level-2 Pellis interference coefficient). Given the depth-D
        sequence (ck), define level-2 interference coefficients

  gammak = ck - (ck-1 ck+1)/(ck-1 + ck+1 + 1), k = 1, …, D-1,

  with boundary conditions gamma0 = c0 and gammaD = cD. The quantity deltak = ck -
  gammak measures the interference correction at level k.

  The resulting tree of correction terms at depth L is the Pellis interference tree P(x, D,
  L), a binary tree of height L with D leaves at each depth. The analogy with a
  multiresolution wavelet decomposition is exact in spirit: the coarse coefficients c0, c1
  carry the leading structure, while the interference corrections deltak at higher depths
  carry fine-grained detail.

        Remark (wavelet analogy). The Pellis interference structure parallels the
        multiresolution analysis of Mallat (1989), with $\varphi$-k playing the role of dilation
        factors and gammak playing the role of scaling coefficients. The key distinction is
        that Mallat's framework operates in L2(R) with orthonormal bases, while the Pellis
        interference tree operates in R>0 with a non-orthogonal, greedy decomposition.
        No claim is made that the two frameworks are equivalent.

  4.3 Convergence and Termination Properties

        Theorem 4.4 (Exponential convergence; Theorem). For any x $\in$ (0, $\in$fty),

        Proof. Follows from Proposition 4.2 by induction on D: the remainder RD $\in$ [0,
        $\varphi$-(D+1)) at each step, so the sum converges exponentially in D. □

  Since $\varphi$-1 approx 0.618, each additional depth level multiplies the approximation error
  by at most $\varphi$-1 approx 0.618, corresponding to a 0.694-bit reduction per step ---
  consistent with the description-length analysis of Section 2.6.

        Corollary 4.5 (Lucas ring termination; Theorem). If x $\in$ Z[$\varphi$] --- i.e., x = a +
        b$\varphi$ for a, b $\in$ Z --- then the standard Pellis expansion terminates in finitely many
        steps with zero remainder.

        Proof. Each step of the greedy algorithm subtracts a Z[$\varphi$] element ck $\varphi$-k from the
        current remainder, producing a new remainder in Z[$\varphi$] cap [0, $\varphi$-k). The set Z[$\varphi$]
        cap [0, $\varepsilon$) is empty for sufficiently small $\varepsilon$ > 0, so the algorithm terminates. An
        explicit bound is that termination occurs in at most 2max(|a|, |b|) steps
        (established in companion Paper 1 of this trilogy, Theorem 3.2). □

  Sanctioned seed verification. For each sanctioned seed F17 = 1597, …, F21 = 10946,
  L7 = 29, L8 = 47: each is an integer, hence in Z[$\varphi$] with b = 0. By Corollary 4.5,
  Algorithm 4.1 terminates at depth 0 with c0 = x and zero remainder. This is a
  regression test for the reference implementation (Appendix B, Test 2; Section 13.5, Test
  P3-T3).

  4.4 Relation to $\beta$-Expansions and Continued Fractions
  The Pellis expansion relates to three classical constructions:

  1. Rényi $\beta$-expansion (Rényi 1957): base-$\varphi$ $\beta$-expansion, digits $\in$ {0,1}, no consecutive
  1s (so the expansion of $\varphi$ in base $\varphi$ is 1.0, and that of $\varphi$2 = $\varphi$ + 1 is 10.01). Pellis relaxes
  the digit constraint and adds hierarchical interference.

  2. Parry characterisation (Parry 1960): characterises bases beta for which the
  expansion of 1 is finite or periodic (Parry numbers). $\varphi$ is a Parry number since its
  beta-expansion of 1 is finite. The Boyd framework extends this to Salem/Pisot number
  classification. The Pellis expansion inherits the Parry-number properties of $\varphi$ but
  relaxes the digit constraint.

  3. Continued fraction expansion: best rational approximations; Pellis gives best
  $\varphi$-polynomial approximations, as developed in companion Paper 2 of this trilogy.

  5. The Projection/Coarse-Graining Map

\subsection{References}

- Vasilev, D. \textit{Flos Aureus Monograph: Trinity Constants}. DOI \href{https://zenodo.org/doi/10.5281/zenodo.19227877}{10.5281/zenodo.19227877}.
- CODATA 2018 recommended values: NIST \href{https://physics.nist.gov/cuu/Constants/}{https://physics.nist.gov/cuu/Constants/}.
- Heyrovska, R. Golden-angle FSC publications --- Heyrovský Institute of Physical Chemistry.
